% =====================================================================
% Capítulo: Pruebas del Sistema
% Fuentes: apps/web/supabase/tests/database (pgTAP), *.test.ts(x) (Vitest),
% apps/e2e/tests (Playwright), .github/workflows/ci.yml,
% docs/tfg/{PROGRESO,BITACORA,TRAZABILIDAD,DECISIONES}.md,
% .claude/skills/rls-review. Cifras: recuento de ficheros con grep y última
% verificación registrada (F3c); unitarias confirmadas con `pnpm test:unit`.
% =====================================================================

% Estilo de los extractos de este capítulo (mismo criterio que en
% Construcción: listings no admite UTF-8 multibyte de forma nativa).
\lstdefinestyle{pymekitpru}{
  basicstyle=\ttfamily\scriptsize,
  breaklines=true,
  frame=single,
  columns=fullflexible,
  keepspaces=true,
  showstringspaces=false,
  literate={á}{{\'a}}1 {é}{{\'e}}1 {í}{{\'i}}1 {ó}{{\'o}}1 {ú}{{\'u}}1
           {ñ}{{\~n}}1 {Á}{{\'A}}1 {É}{{\'E}}1 {Ó}{{\'O}}1 {Ú}{{\'U}}1
}

Este capítulo presenta la estrategia de pruebas de PymeKit y sus
resultados: pruebas unitarias, de integración con la base de datos y la API
del CMS, pruebas de sistema funcionales (E2E por caso de uso) y no
funcionales (seguridad y rendimiento), y pruebas de aceptación. Termina con
las cifras globales y la matriz de trazabilidad entre requisitos y pruebas.

% ---------------------------------------------------------------------
\section{Estrategia de pruebas}
\label{sec:pru-estrategia}

La estrategia sigue el criterio de la pirámide de pruebas: muchas pruebas
rápidas y aisladas en la base, y pocas pruebas de extremo a extremo en la
cima. Su rasgo distintivo es que la base de datos tiene su propio nivel de
pruebas, porque en PymeKit es la frontera de autorización (sección~\ref{sec:dis-seguridad}):
una política RLS mal escrita no se detecta con pruebas de la interfaz, ya que
la interfaz no ofrece el camino para explotarla.

\begin{table}[htbp]
  \centering
  \footnotesize
  \caption{Niveles de prueba, herramientas y qué verifica cada uno.}
  \label{tab:pru-niveles}
  \begin{tabularx}{\textwidth}{@{}>{\raggedright\arraybackslash}p{2.6cm}>{\raggedright\arraybackslash}p{2.4cm}X@{}}
    \toprule
    \textbf{Nivel} & \textbf{Herramienta} & \textbf{Qué verifica} \\
    \midrule
    Unitarias & Vitest & Funciones puras y componentes: filtros, constructor de consultas, formateadores, cálculo de la cuadrícula de paneles, clasificación de errores, componentes de interfaz, coherencia de las claves de traducción \\
    Integración (base de datos) & pgTAP~\cite{pgtap} & Políticas RLS, privilegios, funciones y \textit{triggers}, ejecutados como usuarios reales; aislamiento entre cuentas y entre niveles de personal del CMS \\
    Integración (API del CMS) & Playwright (peticiones HTTP) & Autenticación, permisos y códigos de error de la API montada en la web \\
    Sistema (funcionales) & Playwright~\cite{playwright} & Casos de uso completos en el navegador contra la \textit{build} de pruebas \\
    Sistema (no funcionales) & pgTAP, revisión adversarial, revisión manual & Seguridad (intentos de ataque), rendimiento y usabilidad \\
    Aceptación & Revisión del autor y del tutor & Adecuación a la propuesta \\
    \bottomrule
  \end{tabularx}
\end{table}

Todas las pruebas automáticas forman parte de la integración continua
(sección~\ref{sec:imp-entorno}): las unitarias se ejecutan en cada
\textit{push} a la rama principal y en cada \textit{pull request}; pgTAP y E2E, solo en cada \textit{pull request} hacia la rama
principal. Además, cada cambio se cierra con una lista de verificación
obligatoria (tipos, \textit{lint}, formato, control de marca, pruebas
afectadas, revisión adversarial del \textit{diff} y, si toca la base de
datos, revisión de RLS). Las pruebas E2E se ejecutan siempre contra la
\textit{build} de pruebas y no contra el servidor de desarrollo (B-13).

% ---------------------------------------------------------------------
\section{Pruebas unitarias}
\label{sec:pru-unitarias}

Hay 116 ficheros de pruebas unitarias repartidos en 28 paquetes, con 1.691
pruebas superadas y una omitida. La mayor parte corresponde a los paquetes
del CMS, que concentran la lógica algorítmica: paneles (262 pruebas), núcleo
del explorador de datos (234), constructor de consultas (222), formateadores
(218), núcleo de filtros (181) y componentes del explorador (137).

Algunas pruebas unitarias son, en realidad, pruebas de seguridad. Por
ejemplo, la de \texttt{buildParameterizedStatement} reproduce el valor
malicioso que permitía la inyección SQL heredada de la incidencia B-24
(un nombre que contenía la cadena \texttt{\$2}) y comprueba que ahora viaja
como parámetro enlazado; y las de protección de usuarios verifican que nadie
puede actuar sobre sí mismo ni sobre un super-administrador (B-21). Otras
garantizan propiedades transversales, como que todas las claves de traducción
usadas existen en los mensajes en español.

% ---------------------------------------------------------------------
\section{Pruebas de integración}
\label{sec:pru-integracion}

\subsection{Base de datos: RLS y aislamiento multi-tenant}
\label{sec:pru-pgtap}

La batería pgTAP tiene 70 ficheros (68 de pruebas y 2 de utilidades) con
2.073 pruebas. Cada fichero se ejecuta dentro de una transacción que se
deshace al terminar; crea usuarios de prueba, se \emph{autentica} como cada
uno de ellos (fijando su JWT y el rol \texttt{authenticated} en la sesión,
como hace el servicio de autenticación) y comprueba qué puede ver y hacer.
Se organizan en cuatro grupos:

\begin{itemize}
  \item \textbf{Aislamiento de la plataforma} (cuentas personales y de
        equipo, membresías, invitaciones y sus límites, transferencia de
        propiedad, facturación, notificaciones, almacenamiento,
        super-administrador): un usuario solo ve y modifica lo de sus
        cuentas, y solo con el permiso adecuado.
  \item \textbf{Superficie de privilegios}: \texttt{privileges.test.sql}
        comprueba que ningún rol de la API conserva privilegios residuales
        (como \texttt{TRUNCATE}, que ignora RLS) y \texttt{anon-surface.test.sql}
        fija la superficie exacta del rol anónimo ---solo puede leer el blog
        y no ejecuta ninguna función--- y falla si se amplía (B-39).
  \item \textbf{Control de acceso del CMS} (41 ficheros): permisos de datos,
        sistema y almacenamiento, rangos, delegación, denegaciones,
        auditoría, paneles, vistas guardadas, búsqueda global y el pegamento
        con el super-administrador.
  \item \textbf{Regresión de seguridad}: un fichero por revisión adversarial
        (\texttt{cms-isolation}, \texttt{cms-audit-isolation},
        \texttt{cms-members-isolation}, \texttt{cms-rbac-isolation}\ldots) que
        conserva como prueba cada ataque estudiado.
\end{itemize}

El extracto~\ref{lst:pru-tenant} muestra el patrón básico de una prueba de
aislamiento: tras crear una suscripción para una cuenta, un usuario ajeno no
obtiene ninguna fila.

\begin{lstlisting}[style=pymekitpru, language=SQL, caption={Aislamiento de la facturación entre cuentas (\texttt{team-billing-subscriptions.test.sql}, recortado).}, label={lst:pru-tenant}]
select tests.create_supabase_user('foreigner');
select pymekit.authenticate_as('foreigner');

select is_empty(
  $$ select 1 from subscriptions where id = 'sub_test' $$,
  'The account cannot read the other account subscriptions'
);
\end{lstlisting}

Las pruebas de regresión de las revisiones adversariales ejecutan el ataque
tal como se descubrió. El extracto~\ref{lst:pru-ataque} reproduce la escalada
de la incidencia B-05: antes de la corrección la sentencia afectaba a una
fila; ahora debe fallar con el código de permiso denegado (\texttt{42501}).
La misma prueba comprueba a continuación que la operación legítima sigue
funcionando, para que la corrección no sea simplemente «bloquearlo todo».

\begin{lstlisting}[style=pymekitpru, language=SQL, caption={Regresión de la escalada por \texttt{UPDATE} sin \texttt{WITH CHECK} (\texttt{cms-isolation.test.sql}).}, label={lst:pru-ataque}]
select cms_tests.authenticate_as('iso_mid');

-- El ataque que funcionaba antes de la corrección (UPDATE 1)
select throws_ok(
    $$ update cms.role_permissions
       set permission_id = cms_tests.test_uuid(9304)
       where role_id = cms_tests.test_uuid(9202)
         and permission_id = cms_tests.test_uuid(9303) $$,
    '42501',
    null,
    'E1: no puede colgar de un rol inferior un permiso que no posee'
);
\end{lstlisting}

\subsection{API del CMS}
\label{sec:pru-api}

La API del CMS se prueba de dos formas. Sus servicios y utilidades tienen
pruebas unitarias con el cliente de datos simulado. Y la API montada en la
web se prueba con peticiones HTTP reales desde Playwright
(\texttt{cms-api.spec.ts} y las comprobaciones de API de los demás ficheros
del CMS): sin sesión responde 401; con sesión pero sin acceso al CMS, o sin
segundo factor, o con la cuenta del CMS desactivada, 403; una petición con
campos no permitidos (por ejemplo, metadatos de la aplicación al crear un
usuario) se rechaza con 400; y cada rechazo se comprueba por su código de
error estable, no solo por el estado HTTP, porque un 403 no demuestra nada si
no se sabe por qué se produjo (B-26).

% ---------------------------------------------------------------------
\section{Pruebas de sistema}
\label{sec:pru-sistema}

\subsection{Pruebas funcionales por caso de uso}
\label{sec:pru-funcionales}

Las pruebas E2E se organizan en 33 ficheros (unas 164 pruebas declaradas),
19 de ellos del CMS, con 99 pruebas. Se ejecutan en español, el único idioma
de la interfaz, y un fichero específico comprueba que no queda ningún texto
en inglés visible. La tabla~\ref{tab:pru-cu} relaciona cada caso de uso con
las pruebas que lo cubren.

\begin{table}[htbp]
  \centering
  \footnotesize
  \caption{Cobertura E2E de los casos de uso.}
  \label{tab:pru-cu}
  \begin{tabularx}{\textwidth}{@{}l>{\raggedright\arraybackslash}p{5.0cm}>{\raggedright\arraybackslash}X@{}}
    \toprule
    \textbf{CU} & \textbf{Ficheros de prueba} & \textbf{Estado} \\
    \midrule
    CU-01 & \texttt{authentication/*} (registro, acceso, recuperación de contraseña, redirecciones, sesión caducada), \texttt{account} & Cubierto \\
    CU-02 & \texttt{team-accounts/*}, \texttt{invitations} (incluida la aceptación con MFA) & Cubierto con los roles actuales; roles orientados a pymes pendientes (F5) \\
    CU-03 & \texttt{team-billing}, \texttt{user-billing} & Escritos, pero desactivados mientras no se configuren las claves de la pasarela en modo de pruebas (F5/F6) \\
    CU-04 & pgTAP de facturación y permisos (capa de datos) & Parcial: falta la prueba E2E de la interfaz y de la \textit{server function} (F6) \\
    CU-05 & \texttt{admin} (404 sin permiso, MFA obligatorio, bloquear, reactivar, borrar) & Parcial: la prueba de suplantación está desactivada \\
    CU-06 & \texttt{cms-data-explorer*}, \texttt{cms-audit-logs}, \texttt{cms-settings-*}, \texttt{cms-ui}, \texttt{cms-sidebar-areas}, \texttt{cms-field-help}, \texttt{cms-page-summaries} & Cubierto \\
    CU-07 & --- & Pendiente: evaluación de la reutilización (F6, P-01) \\
    CU-08 & \texttt{blog} (lectura anónima, borradores invisibles, Markdown seguro, publicación desde el CMS, \textit{sitemap}) & Cubierto \\
    \bottomrule
  \end{tabularx}
\end{table}

Las pruebas del CMS verifican sistemáticamente los dos lados de cada
permiso: el personal de soporte, con un rol limitado, \emph{no} ve en la
barra lateral ni puede abrir por URL las tablas sin permiso, y el
super-administrador sí. Muchas incluyen además pruebas negativas
específicas derivadas de incidencias: la guarda de clave en la edición
(B-19), un HTML disfrazado de imagen en la subida de ficheros (B-25), el
\textit{widget} de una tabla no legible («sin acceso», B-50) o la consola con
una sesión sin segundo factor (B-14).

\subsection{Pruebas no funcionales: seguridad}
\label{sec:pru-seguridad}

La seguridad se ha verificado con revisiones adversariales en dos niveles:
una revisión del \textit{diff} orientada a la búsqueda de fallos al terminar
cada funcionalidad, y una revisión específica de RLS (descrita en la
sección~\ref{sec:imp-agentes}) en cada cambio relevante de políticas,
privilegios o funciones de privilegios elevados. La tabla~\ref{tab:pru-rls}
resume las revisiones de RLS realizadas.

\begin{table}[htbp]
  \centering
  \footnotesize
  \caption{Revisiones adversariales de RLS y su resultado.}
  \label{tab:pru-rls}
  \begin{tabularx}{\textwidth}{@{}l>{\raggedright\arraybackslash}p{3.4cm}X@{}}
    \toprule
    \textbf{Fase} & \textbf{Alcance} & \textbf{Resultado} \\
    \midrule
    F2.1 & Esquema \texttt{cms} completo tras su incorporación & Dos brechas heredadas confirmadas (B-05, B-08), cinco endurecimientos y cuatro funciones desalineadas (ADR-015). La refutación independiente encontró B-08 \\
    F2.6 & Auditoría y búsqueda global & Aislado; cuatro debilidades menores endurecidas (B-31 a B-34); +31 pruebas \\
    F2.6b & Blog y apertura del esquema \texttt{public} al rol anónimo & Aislado; riesgo latente de privilegios por defecto documentado y vigilado por prueba (B-39) \\
    F2.7a & Miembros del CMS y autor de la auditoría & Aislado; +68 pruebas; una inyección SQL latente de origen propio (B-45) \\
    F2.7b & RBAC del CMS con administradores delegados de distintos rangos & Brecha heredada de escape de denegaciones mediante delegación, reproducida por el refutador y corregida (B-47) \\
    \bottomrule
  \end{tabularx}
\end{table}

La tabla~\ref{tab:pru-hallazgos} recoge los hallazgos de seguridad más
relevantes de todo el proyecto. De las 29 incidencias de seguridad
registradas, 25 procedían del código de partida, 2 del entorno y 2 se
introdujeron durante el proyecto. Todas están corregidas y cubiertas por
pruebas de regresión.

\begin{longtable}{@{}l>{\raggedright\arraybackslash}p{6.6cm}lll@{}}
  \caption{Principales hallazgos de seguridad corregidos (bitácora del proyecto).}\label{tab:pru-hallazgos}\\
  \toprule
  \textbf{ID} & \textbf{Fallo} & \textbf{Origen} & \textbf{Grav.} & \textbf{Detección} \\
  \midrule
  \endfirsthead
  \toprule
  \textbf{ID} & \textbf{Fallo} & \textbf{Origen} & \textbf{Grav.} & \textbf{Detección} \\
  \midrule
  \endhead
  \bottomrule
  \endfoot
  B-04 & Base de datos local accesible desde Internet & entorno & Alta & Diagnóstico de red \\
  B-05 & Escalada en el RBAC del CMS: \texttt{UPDATE} sin \texttt{WITH CHECK} & heredado & Alta & Revisión de RLS \\
  B-06 & La comprobación de MFA del CMS fallaba en abierto & heredado & Alta & Prueba nueva \\
  B-07 & Lectura de los usuarios de autenticación (con \textit{hashes}) por una función de consulta & heredado & Alta & Revisión de RLS \\
  B-08 & Lectura del catálogo de roles de PostgreSQL & heredado & Alta & Refutación \\
  B-14 & La consola se cargaba sin segundo factor & propio & Media & Prueba manual \\
  B-19 & Edición y borrado de varias filas con una condición que no era clave & heredado & Alta & Revisión del \textit{diff} \\
  B-21 & La gestión de usuarios solo comprobaba el acceso al CMS, no el permiso ni el objetivo & heredado & Alta & Revisión de seguridad \\
  B-24 & Inyección SQL por sustitución manual de parámetros & heredado & Alta & Revisión de informes \\
  B-25 & Subidas con el tipo de contenido elegido por el cliente & heredado & Media & Revisión de seguridad \\
  B-27 & Entradas de auditoría falsificables por el personal & heredado & Alta & Revisión de seguridad \\
  B-28 & La búsqueda global podía buscar en los usuarios de autenticación & heredado & Alta & Revisión de seguridad \\
  B-42 & Un miembro podía cambiar sus propios roles & heredado & Alta & Revisión de seguridad \\
  B-44 & Cualquiera con permiso de ajustes podía desactivar el MFA obligatorio & heredado & Media & Revisión de seguridad \\
  B-45 & Función interna de borrado ejecutable directamente (inyección latente) & propio & Alta & Revisión de RLS \\
  B-47 & Escape de denegaciones explícitas delegando en una cuenta títere & heredado & Alta & Revisión de RLS \\
  B-50 & \textit{Widgets} sobre tablas no legibles y apropiación de paneles compartidos & heredado & Media & Revisión del \textit{diff} \\
\end{longtable}

De estos resultados se extraen tres conclusiones sobre el proceso. Primera,
la demostración con pruebas ejecutables y la refutación independiente
encuentran fallos que la lectura estática da por buenos (B-08, B-47).
Segunda, las revisiones se complementan: B-45, introducido en F2.1, no lo vio
la revisión de esa fase y lo encontró la de F2.7a, que buscaba otra cosa.
Tercera, la defensa en profundidad funciona: en B-14 falló la guarda de la
interfaz, pero no se expuso ningún dato porque la API y la base de datos
seguían exigiendo el segundo factor.

\subsection{Pruebas no funcionales: rendimiento}
\label{sec:pru-rendimiento}

No se han realizado pruebas de carga formales; el requisito RNF-09 pide un
rendimiento \emph{razonable} y se verifica por diseño (carga en el servidor,
caché de consultas, paginación en el servidor con recuento aproximado en
tablas grandes y tiempo máximo de 5~s en la búsqueda global, B-32) y con
pruebas puntuales. La incidencia más ilustrativa (B-55) surgió de forma
espontánea: la base de datos de pruebas, que acumula cuentas en cada
ejecución completa de los E2E, llegó a unas 200 cuentas y la lista de
miembros de un equipo empezó a superar el tiempo máximo de sentencia (8~s),
con consultas de más de 6~s. La causa era una función volátil
(\texttt{active\_account\_id()}) en la cláusula \texttt{WHERE}, que PostgreSQL
evaluaba una vez por cada fila de \texttt{accounts}. Envolverla en
\texttt{(select \ldots)} hace que se evalúe una sola vez como
\textit{InitPlan} y la consulta volvió a ser inmediata. El mismo patrón se
aplica de forma sistemática a \texttt{auth.uid()} en las políticas RLS
(extracto~\ref{lst:imp-rls}).
\todo{Si se quiere reforzar RNF-09, medir tiempos de respuesta de las
pantallas principales con un volumen de datos representativo (por ejemplo,
10.000 cuentas) y añadir la tabla de resultados.}

\subsection{Pruebas no funcionales: usabilidad y accesibilidad}
\label{sec:pru-usabilidad}

La usabilidad y la accesibilidad (RNF-08) se han verificado mediante revisión
manual del autor, que dio lugar a varios rediseños (consola por áreas,
nombres en español, ayuda contextual y resúmenes, ADR-020), y mediante
pruebas E2E que, en buena parte, localizan los elementos por su rol y su
nombre accesible, de modo que fallan si un control pierde su etiqueta; la
prueba de solo español revisa también los atributos accesibles.
\todo{Añadir, si se realiza, una evaluación de accesibilidad automatizada
(por ejemplo, con una herramienta de auditoría de páginas) o con usuarios.}

% ---------------------------------------------------------------------
\section{Pruebas de aceptación}
\label{sec:pru-aceptacion}

Las pruebas de aceptación validan que el sistema satisface la propuesta del
TFG desde el punto de vista de su destinatario. Durante el desarrollo, el
autor ha actuado como cliente: probó a mano cada incremento (por ejemplo,
CU-01 al cierre de F1) y de esas pruebas surgieron incidencias (B-14, B-40,
B-52) y cambios de diseño (ADR-017, ADR-020, ADR-021).
\todo{Completar tras F6: sesión de aceptación con el tutor recorriendo CU-01
a CU-08 con un guion (precondición, pasos, resultado esperado y obtenido),
resultado de la evaluación de la reutilización (CU-07, P-01) y conformidad
final.}

% ---------------------------------------------------------------------
\section{Resultados}
\label{sec:pru-resultados}

La tabla~\ref{tab:pru-cifras} resume el estado de las pruebas en la última
verificación completa (cierre de F3c).

\begin{table}[htbp]
  \centering
  \footnotesize
  \caption{Resultados de las pruebas automáticas (cierre de F3c).}
  \label{tab:pru-cifras}
  \begin{tabularx}{\textwidth}{@{}>{\raggedright\arraybackslash}p{3.4cm}rr>{\raggedright\arraybackslash}X@{}}
    \toprule
    \textbf{Tipo} & \textbf{Ficheros} & \textbf{Pruebas} & \textbf{Resultado} \\
    \midrule
    Unitarias (Vitest) & 116 & 1.692 & 1.691 superadas, 1 omitida \\
    Base de datos (pgTAP) & 70 & 2.073 & Todas superadas \\
    E2E (Playwright) & 33 & $\approx$164 & 156 superadas; el resto, condicionadas a configuración (facturación, gestión de equipos y suplantación) \\
    Control de marca & --- & --- & 0 referencias \\
    \bottomrule
  \end{tabularx}
\end{table}

La evolución de las cifras a lo largo de los incrementos se recoge en la
tabla~\ref{tab:imp-fases}: la batería pgTAP pasó de 396 pruebas heredadas a
2.073, y la E2E de 43 a 156, en buena parte por las pruebas de regresión de
seguridad. Quedan algunas pruebas E2E inestables bajo ejecución paralela, que
pasan al reintentarlas o en serie y cuya estabilización está prevista antes
de la validación final.

La tabla~\ref{tab:pru-trazabilidad} cierra el capítulo con la trazabilidad
entre cada requisito y su evidencia de verificación.

\begin{longtable}{@{}l>{\raggedright\arraybackslash}p{8.4cm}>{\raggedright\arraybackslash}p{3.6cm}@{}}
  \caption{Trazabilidad entre requisitos y pruebas (cierre de F3c).}\label{tab:pru-trazabilidad}\\
  \toprule
  \textbf{Req.} & \textbf{Evidencia} & \textbf{Estado} \\
  \midrule
  \endfirsthead
  \toprule
  \textbf{Req.} & \textbf{Evidencia} & \textbf{Estado} \\
  \midrule
  \endhead
  \bottomrule
  \endfoot
  RF-01 & E2E \texttt{authentication}, \texttt{account}, \texttt{blog}; pgTAP \texttt{blog}, \texttt{anon-surface} & Verificado \\
  RF-02 & Unitarias; E2E de todos los flujos & Verificado \\
  RF-03 & E2E \texttt{authentication} & Verificado \\
  RF-04 & pgTAP \texttt{super-admin}, \texttt{cms-verify-admin-access}; E2E \texttt{admin}, \texttt{team-invitation-mfa}, \texttt{cms-settings-auth} & MFA verificado; OAuth configurable sin prueba \\
  RF-05 & E2E \texttt{account} & Verificado \\
  RF-06 & pgTAP de cuentas, membresías, invitaciones y permisos; E2E \texttt{team-accounts}, \texttt{invitations} & Parcial (roles para pymes en F5) \\
  RF-07 & pgTAP de facturación; E2E de facturación desactivados & Parcial (F5/F6) \\
  RF-08 & E2E \texttt{admin}, \texttt{cms-ui} & Parcial (suplantación sin E2E) \\
  RF-09 & pgTAP del CMS (41 ficheros); unitarias del CMS; E2E del CMS (19 ficheros) & Verificado \\
  RF-10 & pgTAP de auditoría e integridad; E2E \texttt{cms-audit-logs} & Verificado \\
  RF-11 & pgTAP de paneles; unitarias de paneles; E2E \texttt{cms-dashboards} & Verificado \\
  RF-12 & pgTAP \texttt{notifications} & Parcial \\
  RF-13 & Unitaria de coherencia de claves; E2E \texttt{i18n-solo-espanol} & Verificado \\
  RNF-01 & Configuración tipada; evaluación de la reutilización & Pendiente (F6) \\
  RNF-02 & pgTAP; revisiones adversariales (tabla~\ref{tab:pru-rls}) & Verificado \\
  RNF-03 & pgTAP de aislamiento de la plataforma y del CMS & Verificado \\
  RNF-04 & Estructura del monorepo; comprobación de tipos y dependencias & Verificado \\
  RNF-05 & CI (tipos, \textit{lint}, formato, control de marca) & Parcial (comentarios en F4) \\
  RNF-06 & CI con las tres suites & Parcial (completar en F6) \\
  RNF-07 & Despliegue de prueba y manual & Pendiente (F7) \\
  RNF-08 & Revisión manual; E2E por rol y nombre accesible & Parcial \\
  RNF-09 & Revisión del diseño; B-55 & Parcial \\
\end{longtable}
\todo{Actualizar la tabla~\ref{tab:pru-trazabilidad} y las cifras de la
tabla~\ref{tab:pru-cifras} tras F4--F7.}
